\subsubsection*{1.2.7.}
设 $D$ 是圆盘，$0\in D$，令 $D^*=D-\{0\}$，设 $f:X\longrightarrow D$ 为连续映射，并令 $X_0=f^{-1}(0)$。我们将定义函子
\[
\Psi:(\text{层，底空间为 }X)
\longrightarrow(\text{层，底空间为 }X_0\times\mathcal D).
\]
设 $\widetilde D^*$ 是 $D^*$ 的万有覆盖，$\widetilde D$ 如 1.1.9，并令
\[
\overline X=X\times_D\widetilde D,
\qquad
X^*=X-X_0=X\times_DD^*,
\qquad
\overline X^*=\overline X-X_0=X\times_D\widetilde D^*.
\]
有交换图
\begin{equation}\tag{1.2.7.1}
\begin{tikzcd}
X_0 \arrow[r, hook, "\overline i"] \arrow[d, equal]
  & \overline X \arrow[d, "p"]
  & \overline X^* \arrow[l, hook, "\overline j"']
      \arrow[d, "p'"] \\
X_0 \arrow[r, hook, "i"]
  & X
  & X^* \arrow[l, hook, "j"']
\end{tikzcd}
\end{equation}
它在图
\begin{equation}\tag{1.2.7.2}
\begin{tikzcd}
\{0\} \arrow[r, hook, "\overline i"] \arrow[d, equal]
  & \widetilde D \arrow[d, "p"]
  & \widetilde D^* \arrow[l, hook, "\overline j"']
      \arrow[d, "p'"] \\
\{0\} \arrow[r, hook, "i"]
  & D
  & D^* \arrow[l, hook, "j"']
\end{tikzcd}
\end{equation}
之上是笛卡儿的。
\noindent\emph{编校注。} (1.2.7.1) 中原印本右上方的节点缺少上划线。前面的定义及映射 $p'$ 要求该节点为 $\overline X^*$，此处采用这一记号。

对 $X$ 上的层 $F$，定义（采用 1.2.4(c) 的记号）
\[
\Psi(F)=
\bigl(i^*F,\ \overline i^*\overline j_*((jp')^*F),\ \alpha\bigr),
\]
其中
\[
\begin{split}
\alpha:i^*F=\overline i^*p^*F
&\longrightarrow
\overline i^*\overline j_*((jp')^*F)\\
&=\overline i^*\overline j_*\overline j^*p^*F
\end{split}
\]
是在伴随映射 $\mathrm{Id}\longrightarrow\overline j_*\overline j^*$ 上施加 $\overline i^*$ 而得。

还定义函子
\[
\Psi_\eta:(\text{层，底空间为 }X^*)
\longrightarrow(\text{层，底空间为 }X_0\times\mathcal D^*)
\]
令
\[
\Psi_\eta(F)=\overline i^*\overline j_*({p'}^*F).
\]

\subsubsection*{1.2.8.}
$\Psi$ 的函子性质与 XIII, 1.3.5--1.3.9 中研究的对应函子相同。

